\section{实验}

\subsection{实验设置}

实验基于东海274个开发集窗口（2018至2022年，用于标定底图误差、绘制相图及训练历史基线）与148个测试集窗口（2023至2026年，用于主实验与消融检验）开展，参数汇总于表\ref{tab:settings}。测试涵盖慢计算机（5次/分，$+20\text{ W}$）、中等计算机（20次/分，$+20\text{ W}$）、快计算机（100次/分，$+20\text{ W}$）与省电加速器（20次/分，$+5\text{ W}$）四种硬件，以及“当天卫星图 + 小尺度变化”（主设置，$s_{\text{sat}}=0.15$）和“当天卫星图”（$s_{\text{sat}}=0$）两种真值模型。同一任务上所有控制器共享相同的真实海洋、起点与传感器噪声序列，且随机种子按任务固定（慢计算机独立重跑全部148个测试任务预演所得遗憾与首次运行逐任务一致）；配对差值的95\%置信区间由按赤潮周整组重抽样的自助法给出。

\begin{table}[htbp]
  \centering
  \caption{东海赤潮应急巡测实验的核心参数与评价设置。}
  \label{tab:settings}
  \begin{tabular}{@{}p{0.17\linewidth}p{0.79\linewidth}@{}}
\toprule
项目 & 设置 \\
\midrule
卫星数据 & 美国国家海洋和大气管理局海岸观测网多星融合叶绿素 a 产品（经验正交函数插补去云），2 km、逐日，2018-01-01 至 2026-10-23 共 3197 天；喷气推进实验室多尺度融合海表温度，0.01° \\
任务窗口 & 每年 4--6 月浙江外海陆架（26.5°--31.5°N，121.5°--125°E）局部峰值超过 5 mg/m$^3$ 处，80 km $\times$ 80 km（40 $\times$ 40 个 2 km 网格）；开发 274 个（2018--2022 年），测试 148 个（2023--2026 年） \\
真实海洋 & 任务当天卫星图叠加卫星分辨不出的小尺度变化（标准差 0.15，相关尺度 6 km，主设置）；或任务当天卫星图本身 \\
航行器 & 航速 1.5 m/s，电池 1 kWh；平台 15 W，推进 30 W，高功耗芯片 +20 W（省电加速器 +5 W）；任务 24 h，耗尽电量即停 \\
决策周期 & 20 min；走一段约 1.8 km、采样约 2 次；停下漂流采样 5 次；荧光计单次误差 0.1 \\
想象海洋 & 48 个（32 个用于搜索航线，16 个留出），由 3 天前卫星图与数据标定的误差生成，空间相关长度 10 km，随每个样本同化更新 \\
规划器 & 交叉熵搜索经两个航路点的折线航线，跨周期热启动；低功耗处理器每分钟 5、20 或 100 次推演（条件扫描 5--200 次），高功耗芯片为其 50 倍 \\
出发前预演 & 4 个留出海洋作假想真值 $\times$ 5 种用时策略，各完整模拟 24 h，选平均假想遗憾最小者 \\
评价 & 时间平均遗憾（对数浓度，越低越好）；配对差值的 95\% 置信区间由按赤潮周整组的自助重抽样得到 \\
\bottomrule
\end{tabular}

\end{table}

我们将“出发前预演（本文）”与五类控制方式对比：（1）五种固定策略及“五种策略随机选”、“三种不停策略随机选”；（2）需要历史任务的方法，包括“调好的固定策略”（在120个开发集任务上为每种硬件和真值模型选出平均遗憾最低的固定策略）与“按历史任务学习的选择器”（参照学习型元推理\cite{magnino2026metareasoning}，在同批开发集任务上用旧图最大值、中位数、极差、标准差、起点至旧图峰值距离、可航行占比、起点旧图浓度及近峰值网格占比共8项特征训练梯度提升回归预测各策略遗憾并择优）；（3）在线元推理，包括经典“单步元推理”\cite{russell1991right}与“单步元推理（留出估值）”；（4）不含多步航线规划的巡测规则（全程巡航并用相同局地化集合滤波更新地图），包括间距$16\text{ km}$平行扫测的“割草机式全覆盖”\cite{galceran2013coverage}、“贪心：驶向估计最浓处”（驶向集合均值最高未访格）、“上置信界采样”（均值加2倍标准差最高未访格）\cite{srinivas2012gpucb}与“汤普森采样”（每周期随机抽取一个想象海洋并驶向其最浓未访格）\cite{russo2018thompson}；（5）事后挑出最佳固定策略的“每次任务事后最优（参考）”。

\subsection{卫星大数据的作用与局限}

八年卫星叶绿素大数据能够标定滞后底图在不同季节与滞后天数下的误差分布，但深度时空预测模型无法改善1至3天的底图外推精度，小尺度结构必须依靠现场采样获取（图\ref{fig:data}）。全海域3197天数据上相隔3天的对数叶绿素变化标准差在春季赤潮高发期最高（中位数0.127），夏秋季次之，冬季最低（0.081），夏季浙闽沿岸与长江冲淡水锋面形成0.25至0.3的强变率带（图\ref{fig:data}a至图\ref{fig:data}d）。在274个开发集窗口内，滞后1、2、3天旧图的误差标准差依次增至$0.125$、$0.176$与$0.210$（均值偏差$+0.016$、$+0.032$、$+0.045$，相关尺度约5格即$10\text{ km}$），而同日卫星图与现场海水还存在约0.15的对数差异\cite{ye2021cocts}（图\ref{fig:data}e）。

\begin{figure}[htbp]
  \centering
  \includegraphics[width=\linewidth]{fig_data}
  \caption{八年东海卫星叶绿素大数据提供的先验误差标定及其外推预测边界。（a）至（d）2018至2026年全部3197天对数叶绿素3天变化标准差的四季空间分布（中位数春季0.127、夏季0.107、秋季0.103、冬季0.081），春季子图标出了274个开发集窗口与148个测试集窗口的位置。（e）开发集窗口内滞后1、2、3天旧卫星图的误差标准差（0.125、0.176、0.210）及同日卫星图相对于现场实测水体的约0.15偏差。（f）2023至2026年测试集上深度时空预测模型与直接沿用最新卫星图在提前1、2、3天预报上的均方根误差对比（0.1229对0.1226、0.1439对0.1439、0.1487对0.1474）。数据能够可靠标定旧图的不确定性幅值，但无法替代现场观测。}
  \label{fig:data}
\end{figure}

为检验出发前能否用深度学习外推旧图以减小滞后误差，我们训练了550万参数的联合嵌入预测架构模型\cite{assran2025vjepa2}，以过去3天叶绿素和海表温度异常场为输入（2018至2021年3至6月训练，2022年验证集误差仅比沿用最新底图低1.4\%）。在2023至2026年测试集上，该模型提前1、2、3天预报的均方根误差（集合离散度与误差之比$0.70\text{--}0.93$）与直接沿用最新可用卫星图的持久性基线持平（图\ref{fig:data}f）。这说明了海洋大数据在本问题中的角色：数值外推无法替代现场采样，但多年遥感统计能够为生成想象海洋提供准确的先验不确定性标尺。

\subsection{规划自评的偏差}

随着单次决策推演预算增加，规划器在搜索用海洋上的自评得分因优化器诅咒持续虚高，而在留出海洋上评判则能准确反映真实航线价值的饱和过程（图\ref{fig:curse}）。在80个开发集窗口（每窗口3个种子，3天前旧图，任务前8小时，主真值设置）上，当推演次数从64次增至2000次再扩至200,000次时，真实航线价值在2000次后即进入平台期并收敛至$0.803$，而32个搜索用海洋上的自评得分却一路升至$0.917$，在最大预算下高估达$0.114$；同一批航线在16个留出海洋上的评判均值在全部预算下与真实价值之差均不超过$0.005$（图\ref{fig:curse}a、图\ref{fig:curse}b）。

\begin{figure}[htbp]
  \centering
  \includegraphics[width=\linewidth]{fig_curse}
  \caption{推演预算从64次扩展至200,000次时规划器自评得分的乐观偏差与留出海洋的无偏评判（80个开发集窗口$\times 3$个种子，3天前卫星图，任务前8小时）。（a）所选航线的时间平均最佳对数叶绿素价值：搜索用海洋上的自评得分从0.762持续升至0.917，而真实航线价值（0.749升至0.803）与留出海洋评判（0.745升至0.802）在约2000次推演后一同趋于饱和。（b）估计价值减去真实价值的偏差曲线：使用数据标定不确定性时留出评判偏差保持在0.005以内，将不确定性设小一半或放大至1.5倍虽产生$-0.06$与$+0.07$的整体平移，但仍能准确捕捉增加算力的相对增益。}
  \label{fig:curse}
\end{figure}

当人为将先验标准差$\sigma_3$减半或放大至1.5倍时，留出评判虽产生整体平移，但推演预算从64次增至200,000次时留出海洋估计的价值增量（$+0.059$与$+0.054$）仍与真实价值增量（$+0.054$与$+0.047$）吻合（图\ref{fig:curse}b）。这表明将搜索用海洋与留出海洋分离能够切断优化器对样本噪声的迎合，为衡量计算价值提供可靠的独立标尺。

\subsection{主要结果}

在148个测试任务上，出发前预演在四种硬件与两种真值模型下无需历史任务即达到“调好的固定策略”与“按历史任务学习的选择器”水平，并优于经典“单步元推理”以及除慢计算机主设置下贪心规则（持平）外的全部无规划巡测规则（图\ref{fig:main}，表\ref{tab:main}）。在主设置下，预演在四种硬件上的时间平均遗憾为$0.265\text{--}0.276$，与“调好的固定策略”的配对差值为$-0.003$至$+0.007$（慢计算机$[-0.000, +0.014]$、中等计算机$[-0.009, +0.003]$），与“按历史任务学习的选择器”差值为$-0.004$至$+0.000$，全部95\%置信区间均包含0。

\begin{figure}[htbp]
  \centering
  \includegraphics[width=\linewidth]{fig_main}
  \caption{148个测试任务上各对比控制方式减去“出发前预演（本文）”的时间平均遗憾配对差值及95\%置信区间（正值表示差于本文方法）。四列对应慢计算机、中等计算机、快计算机与省电加速器，实心圆点为“当天卫星图 + 小尺度变化”主设置，空心圆点为“当天卫星图”真值设置。出发前预演在所有硬件与两种真值模型下均匹配需要历史任务的两种方法，并显著优于经典单步元推理及所有巡测规则（仅慢计算机主设置下与贪心规则持平）。}
  \label{fig:main}
\end{figure}

\begin{table}[htbp]
  \centering
  \caption{2023至2026年148个测试任务上四种机载硬件与两种真值模型下各控制方式的时间平均遗憾（越低越好，每列最优可运行控制器加粗，本文方法行加底纹）。}
  \label{tab:main}
  \small\resizebox{\linewidth}{!}{\begin{tabular}{@{}llcccccccc@{}}
\toprule
 & & \multicolumn{4}{c}{真值：当天卫星图 + 小尺度变化} & \multicolumn{4}{c}{真值：当天卫星图} \\
\cmidrule(lr){3-6}\cmidrule(l){7-10}
类别 & 控制方式 & 慢计算机 & 中等计算机 & 快计算机 & 省电加速器 & 慢计算机 & 中等计算机 & 快计算机 & 省电加速器 \\
\midrule
固定策略 & 一直走 & 0.374 & 0.267 & 0.267 & 0.267 & 0.281 & 0.196 & 0.194 & 0.196 \\
 & 一直深想 & 0.273 & 0.269 & 0.269 & 0.267 & 0.195 & 0.194 & 0.193 & \textbf{0.189} \\
 & 三段深想一段 & \textbf{0.270} & 0.269 & \textbf{0.265} & 0.266 & \textbf{0.194} & \textbf{0.190} & \textbf{0.191} & 0.191 \\
 & 三段停下深想一段 & 0.331 & 0.328 & 0.327 & 0.327 & 0.252 & 0.248 & 0.245 & 0.247 \\
 & 三段停下复测一段 & 0.408 & 0.335 & 0.331 & 0.335 & 0.319 & 0.251 & 0.248 & 0.251 \\
\midrule
需要历史任务 & 调好的固定策略$^{a}$ & \textbf{0.270} & 0.269 & 0.267 & 0.267 & \textbf{0.194} & 0.196 & \textbf{0.191} & 0.191 \\
 & 按历史任务学习的选择器 & 0.278 & 0.269 & 0.273 & \textbf{0.266} & 0.200 & 0.197 & \textbf{0.191} & 0.189 \\
\midrule
在线元推理 & 单步元推理 & 0.287 & 0.287 & 0.282 & 0.278 & 0.202 & 0.202 & 0.204 & 0.199 \\
 & 单步元推理（留出估值） & 0.278 & 0.270 & 0.270 & 0.267 & -- & -- & -- & -- \\
\midrule
巡测规则 & 割草机式全覆盖 & 0.424 & 0.424 & 0.424 & 0.424 & 0.328 & 0.328 & 0.328 & 0.328 \\
 & 贪心：驶向估计最浓处 & 0.285 & 0.285 & 0.285 & 0.285 & 0.206 & 0.206 & 0.206 & 0.206 \\
 & 上置信界采样 & 0.299 & 0.299 & 0.299 & 0.299 & 0.224 & 0.224 & 0.224 & 0.224 \\
 & 汤普森采样 & 0.332 & 0.332 & 0.332 & 0.332 & 0.256 & 0.256 & 0.256 & 0.256 \\
\midrule
\rowcolor{violet!8}本文 & 出发前预演 & 0.276 & \textbf{0.265} & 0.268 & \textbf{0.266} & 0.197 & 0.191 & 0.195 & 0.190 \\
\midrule
参考 & 每次任务事后最优 & 0.250 & 0.243 & 0.243 & 0.242 & 0.183 & 0.179 & 0.180 & 0.179 \\
\bottomrule
\end{tabular}
}
  \par\vspace{2pt}\footnotesize 注：$^a$ 调好的固定策略在主设置下依次为三段深想一段、三段深想一段、一直走、一直深想，在“当天卫星图”真值下依次为三段深想一段、一直走、三段深想一段、三段深想一段。
\end{table}

相比之下，主设置下预演比经典“单步元推理”降低遗憾$0.011$（$[-0.019, -0.002]$）、$0.022$（$[-0.030, -0.013]$）、$0.013$（$[-0.020, -0.007]$）与$0.012$（$[-0.019, -0.004]$），置信区间全部低于0。对四种无规划巡测规则，预演除在慢计算机上与贪心规则差值为$-0.009$（$[-0.022, +0.004]$）外，在其余三种硬件上比贪心规则低$0.017\text{--}0.020$，并在所有硬件上比上置信界采样、汤普森采样和割草机式全覆盖分别低$0.023\text{--}0.034$、$0.056\text{--}0.067$和$0.148\text{--}0.159$（区间全部低于0）。

在“当天卫星图”真值下，因无亚网格随机场，各控制器绝对遗憾整体下降，但相对格局一致（表\ref{tab:main}右半部）：预演（$0.191\text{--}0.197$）与调好策略及学习型选择器的差值在$-0.005$至$+0.003$之间（区间均含0），比单步元推理低$0.006\text{--}0.010$，比全部四种巡测规则低$0.010\text{--}0.139$（区间全部低于0）；较优停下策略在两种真值下均比最佳不停策略差$0.054\text{--}0.062$。此外，两种真值下148个测试任务间的遗憾标准差达约$0.13$与$0.11\text{--}0.12$，而三种不停策略间的平均差距（除慢计算机外）不超过$0.01$。

\subsection{机理分析}

经典单步元推理失利源于在80\%至97\%的周期内开启高功耗芯片并在33\%至34\%的周期内停下复测，而“单步元推理（留出估值）”将深想比例压低至30\%至44\%、停下比例降至15\%至18\%（预演深想占26\%至54\%、停下仅$0.5\%\text{--}1.8\%$，图\ref{fig:mech}a、图\ref{fig:mech}b），使四种硬件上的遗憾降至$0.278$、$0.270$、$0.270$与$0.267$（比经典单步元推理低$0.009$、$0.017$、$0.012$与$0.010$，与预演之差在$0.005$以内且区间含0）。为分离两项改动的作用，我们在慢计算机与中等计算机上（相同148个任务与种子）分别只做其中一项改动（图\ref{fig:mech}a、b中的空心标记）。只改为留出估值时，深想时段占比即由80\%与93\%降至29\%与38\%，但停下占比仍为34\%与36\%，遗憾为$0.291$与$0.283$（相对经典方法分别为$+0.004$（$[-0.004, +0.013]$）与$-0.004$（$[-0.011, +0.003]$））；只加两倍标准误规则时，深想占比仍为81\%与93\%，停下占比降至26\%与29\%，遗憾为$0.280$与$0.285$（$-0.007$（$[-0.012, -0.003]$）与$-0.002$（$[-0.005, +0.000]$））；两处同时改动才把停下占比均降至17\%，遗憾降至$0.278$与$0.270$（$-0.009$（$[-0.018, -0.000]$）与$-0.017$（$[-0.025, -0.009]$）），达到预演水平。可见过度深想源于在搜索用海洋上估计计算收益，过多停船源于对微小样本差值的反应，二者须同时纠正：单独任一改动在中等计算机上都未明显降低遗憾，在深想本身有益的慢计算机上也只有保守切换规则单独有效。

\begin{figure}[htbp]
  \centering
  \includegraphics[width=\linewidth]{fig_mech}
  \caption{在线元推理过度规划的成因与出发前预演的硬件自适应机制（主真值设置，148个测试任务）。（a）高功耗芯片开启周期占比与（b）停下复测周期占比：经典单步元推理在80\%至97\%的周期内开启高功耗芯片并在33\%至34\%的周期内停下，而单步元推理（留出估值）与出发前预演大幅减少了深想与停船；慢、中计算机行中的空心方块与空心三角分别为只改为留出估值、只加两倍标准误切换规则的诊断变体。（c）出发前预演在四种硬件上选出各候选策略的任务占比：随基础算力提升，“一直走”占比从1\%升至30\%再升至47\%，而在省电加速器上“一直深想”占比从25\%升至43\%。}
  \label{fig:mech}
\end{figure}

出发前预演选出的策略比例随硬件算力与功耗自适应调整（图\ref{fig:mech}c）：在“一直走”比调好策略差$+0.104$的慢计算机上，预演极少选“一直走”，随基础算力升至中等和快计算机，“一直走”占比阶梯式上升，而在同为20次/分但芯片仅$+5\text{ W}$的省电加速器上，“一直深想”占比明显增加。预演的优势来自剔除当前硬件下的劣质选项（各硬件上的停下策略及慢计算机上的“一直走”），故显著优于“五种策略随机选”（$0.292\text{--}0.331$）和慢计算机上的“三种不停策略随机选”（$0.306$）；而在真实差距不超过$0.004$的其余三种硬件不停策略间，其单任务选择接近随机，与“预演（以搜索过的海洋为真值）”选出同一策略的比例为37\%至41\%，但两者平均遗憾仅差$-0.004$至$+0.004$（区间含0），因为预演以假想真值上的最终任务遗憾而非规划器自评得分比较策略。

\begin{figure}[htbp]
  \centering
  \includegraphics[width=\linewidth]{fig_example}
  \caption{一次典型测试任务（2023年4月15日，慢计算机）中“一直走”与“三段深想一段”的航迹及实时遗憾对比。（a）航行器出发时携带的3天前卫星图及东侧边缘入水起点（白色三角形）。（b）真实海洋叶绿素分布（西南角五角星为真实最浓处）以及“一直走”（蓝线）与“三段深想一段”（橙线，加粗段为开启高功耗芯片深想）的24小时航迹。（c）24小时内报告位置的实时遗憾曲线：“一直走”因搜索不足在窗口中部绕圈，遗憾停留在0.4附近（均值0.43），“三段深想一段”横穿窗口并在约19小时后抵达真实最浓处（均值降至0.31）。}
  \label{fig:example}
\end{figure}

在2023年4月15日测试任务（慢计算机，按日期顺序首个两策略差距超0.1的任务，图\ref{fig:example}）中，航行器从东侧入水，真实最浓处位于西南角。慢计算机上“一直走”每周期仅后台推演100次，航行至窗口中部次高值区后局部徘徊，时间平均遗憾为$0.43$；而“三段深想一段”在深想周期内执行5000次推演，规划出的航线横穿窗口指向西南角，于约19小时后抵达真实最浓处，将平均遗憾降至$0.31$。

\subsection{消融实验}

在慢计算机与中等计算机上的消融实验（图\ref{fig:ablation}）首先检验了假想真值的数量与来源。将假想真值从4个减为1个或2个时，两台计算机上的遗憾恶化$+0.006$至$+0.013$；增至8个仅在慢计算机上微降$-0.007$（$[-0.012, -0.001]$）而中等计算机变化$+0.003$，故4个真值是兼顾精度与耗时的合理配置。将假想真值改从搜索用海洋中抽取时，遗憾变化仅为$-0.004$与$+0.003$（区间含0），印证以假想真值上的任务结果评判策略可规避自评过拟合。

\begin{figure}[htbp]
  \centering
  \includegraphics[width=\linewidth]{fig_ablation}
  \caption{慢计算机与中等计算机（主真值设置，148个测试任务）上的消融实验结果（各变体减去完整“出发前预演（本文）”的时间平均遗憾配对差值及95\%置信区间）。减少假想真值数量、将先验不确定性错误放大至2倍或取消预演改为随机选择均导致性能劣化，而在任务尺度预演、在线周期级留出估值与需要历史任务的方法之间表现相当。}
  \label{fig:ablation}
\end{figure}

第二组消融检验先验不确定性标定与评估尺度。将先验标准差减半时遗憾变化不大（$+0.002$与$+0.004$），但放大至标定值的2倍时，慢、中计算机遗憾恶化$+0.011$（$[+0.001, +0.020]$）与$+0.020$（$[+0.011, +0.030]$），此时过大噪声使所有固定策略退化（中等计算机不停策略恶化至$0.275\text{--}0.283$），预演（$0.286$）亦不再匹配最佳固定策略（$0.275$）。将全任务预演替换为在线“单步元推理（留出估值）”时差值仅为$+0.002$与$+0.005$（区间含0），说明整场预演与在线留出估值两种用法效果相当，而取消预演改为五种策略随机选则损失$+0.055$与$+0.028$。

\subsection{适用范围}

在低功耗处理器速度与真实海洋偏离卫星图程度构成的二维相图（图\ref{fig:phase}）上，深想仅在基础算力低于航线价值饱和阈值时明显获益，而停下复测在全部格点均造成损失。在5次/分处理器上（后台每周期仅推演100次，远低于约2000次的饱和点），当底图偏离度为0.05、0.15与0.30时，较优深想策略比“一直走”分别降低遗憾$0.086$、$0.109$与$0.138$，底图越不准深想收益越大；10次/分上仅偏离0.15一格有微弱收益（$-0.017$），从20次/分起均无明显收益。在偏离度为0的第四行（当天卫星图本身为真值），5次/分计算机上深想比“一直走”降低遗憾$0.080$（$[-0.119, -0.043]$），10次/分及以上无明显收益，而该行每一格中较优停下策略均比最佳不停策略差$0.061\text{--}0.076$。

\begin{figure}[htbp]
  \centering
  \includegraphics[width=\linewidth]{fig_phase}
  \caption{低功耗处理器速度（每分钟5至200次推演）与真实海洋偏离卫星图程度（0、0.05、0.15、0.30）二维网格上的策略收益相图（每格60个开发集任务）。（a）最佳深想策略减去“一直走”的时间平均遗憾差值（加框格点表示95\%置信区间低于0，即深想显著有益）：深想仅在基础算力过慢（5次/分）时显著获益，且底图偏离越大收益越强。（b）较好停下策略减去最佳不停策略的遗憾差值：所有格点差值均为显著正值（$+0.061$至$+0.118$），表明停下原位复测在任何算力与底图误差下均得不偿失。}
  \label{fig:phase}
\end{figure}

跨越全部24个相图格点，较优停下策略始终比最佳不停策略高出$0.061\text{--}0.118$的遗憾（图\ref{fig:phase}b）：原位复测5次虽将单点测量误差从$0.1$降至约$0.045$，但巡测瓶颈在于未访水域约$0.26$的先验不确定性，且停船三分之一时间推迟了到达高浓水体的时刻，增大了时间平均遗憾。此外，当改用滞后1天的较新卫星图（$s_1=0.125$，主设置）时，慢计算机上调好策略为“一直深想”（$0.250$），预演取得$0.251$（差值$+0.001$，$[-0.005, +0.011]$），优于单步元推理（$0.260$，差值$-0.009$，$[-0.018, +0.001]$）；中等计算机上调好策略为“一直走”（$0.250$），预演取得$0.243$，比调好策略低$0.007$（$[-0.014, -0.002]$）且比单步元推理（$0.260$）低$0.017$（$[-0.022, -0.013]$），这是因为“一直走”在开发年份最优，而在测试年份上“三段深想一段”（$0.239$）反而是最佳固定策略。

\subsection{局限性}

本研究仍存在若干简化。首先，实验以任务当天卫星叶绿素图（叠加或不叠加标定小尺度场）作为真实海洋而非现场走航剖面，二维随水漂移参考系忽略了水平流剪切与藻类昼夜垂直迁移，双航路点规划器也是对三维流场路径规划的简化替代。其次，预演仅在下水前执行一次而未在任务中途重跑，且4个假想真值主要使策略适配机载硬件特性，尚不足以在平均差距微小的几种不停策略间针对单次任务做出精细区分（主设置下距“每次任务事后最优（参考）”仍差$0.022\text{--}0.027$）。最后，预演中的假想遗憾比真实海洋遗憾整体偏高约$0.05\text{--}0.07$，但该偏移对各策略一致，不改变相对排序。
